\scanpage{34}
% archive sheet number: 33
sont les sections de ramification de $P^1_S$ sur $P^1_S/(Z/3)$. Passant au revêtement galoisien $S'=S(\sqrt{-3})$ de $S$, d'où dans $P^1_{S'}$ \deleted{\unsure{une donnée de|un diviseur de} descente} \inserted{deux sections} par $\overline{S'}$, soient $x'$ et $x''$, qui donnent sur $P^1_{S'}/(Z/3)$ deux sections; prenant pour l'une \deleted{la} \inserted{comme} section unité \unsure{(|(le} $\infty$ et $0$ étant déjà choisis ci-dessus) on trouve que l'autre est $-1$. Cela définit un isomorphisme $(P^1_S/Z/3)\times_S S'\simeq P^1_{S'}$. Ainsi, $P^1_S/(Z/3)$ est can. isomorphe à un fibré principal de groupe $G_{m_S}$, [associé à un faisceau inversible $\Lambda_S$], $\Lambda_S$ étant muni d'une bi-section \inserted{non ramifiée} symétrique non ramifiée, ou (ce qui revient au même) d'une forme quadratique non dégénérée; la \deleted{\unsure{fi|f|fait}} \unsure{classe|de munir} de cette bisection n'étant autre que celle du revêtement $S'=S(\sqrt{-3})$.

Passant à $S'=S(\sqrt{-3})$, on obtient la \unsure{tour|forme} suivante

